\section[Operational Constraints]{Operational Constraints\\and Reachability}
\label{sec:reachability}

The framework's operational picture is almost as important as its vulnerability
picture. The model registry contains 20 aliases, but only seven were executed in
the archived comparison run. \Cref{fig:provider-status} visualizes that
distribution, and \cref{tab:reachability} summarizes it by provider.

\begin{figure}[t]
  \centering
  \includegraphics[width=\columnwidth]{provider_status.pdf}
  \caption{Provider-level reachability for the current registry. Successful rows
  correspond to executed comparison-run models; blocked rows correspond to later
  preflight failures for the remaining aliases.}
  \label{fig:provider-status}
\end{figure}

\begin{table}[t]
  \centering
  \caption{Reachability summary over the 20 configured aliases.}
  \label{tab:reachability}
  \resizebox{\columnwidth}{!}{\begin{tabular}{lrrl}
\toprule
Provider & Registry aliases & Executed or reachable & Dominant blocker \\
\midrule
Anthropic & 5 & 5 & None \\
Google & 1 & 1 & None \\
Mistral & 2 & 2 & None \\
OpenAI & 3 & 3 & None \\
\bottomrule
\end{tabular}
}
\end{table}

The blocker distribution is lopsided but interpretable. All five Anthropic
aliases and both Mistral-hosted aliases were reachable in the archived run. By
contrast, all eight Together-backed aliases failed due to invalid API keys. All
three OpenAI aliases failed due to insufficient quota. The two Google aliases
failed for different reasons: \texttt{gemini-2.0-flash} due to quota exhaustion
and \texttt{gemini-2.5-pro} because the configured preview model identifier had
gone stale and returned a 404.

This is not just infrastructure noise. Operational reachability determines which
claims can be made at all. A registry may look broad on paper, but without valid
credentials, non-expired model identifiers, and adequate quota, that breadth is
not empirically realized. We therefore treat the 13 blocked aliases as part of
the study outcome, not as omitted background conditions.

The reachability story also clarifies a methodological point. The paper's
empirical claims are about the seven executed models. The provider-status table
extends that view by showing what the surrounding registry looked like when the
remaining aliases were probed. It should not be misread as a second comparison
run over all 20 models.
